\documentclass[11pt]{article}

    \usepackage[breakable]{tcolorbox}
    \usepackage{parskip} % Stop auto-indenting (to mimic markdown behaviour)
    
    \usepackage{iftex}
    \ifPDFTeX
    	\usepackage[T1]{fontenc}
    	\usepackage{mathpazo}
    \else
    	\usepackage{fontspec}
    \fi

    % Basic figure setup, for now with no caption control since it's done
    % automatically by Pandoc (which extracts ![](path) syntax from Markdown).
    \usepackage{graphicx}
    % Maintain compatibility with old templates. Remove in nbconvert 6.0
    \let\Oldincludegraphics\includegraphics
    % Ensure that by default, figures have no caption (until we provide a
    % proper Figure object with a Caption API and a way to capture that
    % in the conversion process - todo).
    \usepackage{caption}
    \DeclareCaptionFormat{nocaption}{}
    \captionsetup{format=nocaption,aboveskip=0pt,belowskip=0pt}

    \usepackage[Export]{adjustbox} % Used to constrain images to a maximum size
    \adjustboxset{max size={0.9\linewidth}{0.9\paperheight}}
    \usepackage{float}
    \floatplacement{figure}{H} % forces figures to be placed at the correct location
    \usepackage{xcolor} % Allow colors to be defined
    \usepackage{enumerate} % Needed for markdown enumerations to work
    \usepackage{geometry} % Used to adjust the document margins
    \usepackage{amsmath} % Equations
    \usepackage{amssymb} % Equations
    \usepackage{textcomp} % defines textquotesingle
    % Hack from http://tex.stackexchange.com/a/47451/13684:
    \AtBeginDocument{%
        \def\PYZsq{\textquotesingle}% Upright quotes in Pygmentized code
    }
    \usepackage{upquote} % Upright quotes for verbatim code
    \usepackage{eurosym} % defines \euro
    \usepackage[mathletters]{ucs} % Extended unicode (utf-8) support
    \usepackage{fancyvrb} % verbatim replacement that allows latex
    \usepackage{grffile} % extends the file name processing of package graphics 
                         % to support a larger range
    \makeatletter % fix for grffile with XeLaTeX
    \def\Gread@@xetex#1{%
      \IfFileExists{"\Gin@base".bb}%
      {\Gread@eps{\Gin@base.bb}}%
      {\Gread@@xetex@aux#1}%
    }
    \makeatother

    % The hyperref package gives us a pdf with properly built
    % internal navigation ('pdf bookmarks' for the table of contents,
    % internal cross-reference links, web links for URLs, etc.)
    \usepackage{hyperref}
    % The default LaTeX title has an obnoxious amount of whitespace. By default,
    % titling removes some of it. It also provides customization options.
    \usepackage{titling}
    \usepackage{longtable} % longtable support required by pandoc >1.10
    \usepackage{booktabs}  % table support for pandoc > 1.12.2
    \usepackage[inline]{enumitem} % IRkernel/repr support (it uses the enumerate* environment)
    \usepackage[normalem]{ulem} % ulem is needed to support strikethroughs (\sout)
                                % normalem makes italics be italics, not underlines
    \usepackage{mathrsfs}
    

    
    % Colors for the hyperref package
    \definecolor{urlcolor}{rgb}{0,.145,.698}
    \definecolor{linkcolor}{rgb}{.71,0.21,0.01}
    \definecolor{citecolor}{rgb}{.12,.54,.11}

    % ANSI colors
    \definecolor{ansi-black}{HTML}{3E424D}
    \definecolor{ansi-black-intense}{HTML}{282C36}
    \definecolor{ansi-red}{HTML}{E75C58}
    \definecolor{ansi-red-intense}{HTML}{B22B31}
    \definecolor{ansi-green}{HTML}{00A250}
    \definecolor{ansi-green-intense}{HTML}{007427}
    \definecolor{ansi-yellow}{HTML}{DDB62B}
    \definecolor{ansi-yellow-intense}{HTML}{B27D12}
    \definecolor{ansi-blue}{HTML}{208FFB}
    \definecolor{ansi-blue-intense}{HTML}{0065CA}
    \definecolor{ansi-magenta}{HTML}{D160C4}
    \definecolor{ansi-magenta-intense}{HTML}{A03196}
    \definecolor{ansi-cyan}{HTML}{60C6C8}
    \definecolor{ansi-cyan-intense}{HTML}{258F8F}
    \definecolor{ansi-white}{HTML}{C5C1B4}
    \definecolor{ansi-white-intense}{HTML}{A1A6B2}
    \definecolor{ansi-default-inverse-fg}{HTML}{FFFFFF}
    \definecolor{ansi-default-inverse-bg}{HTML}{000000}

    % commands and environments needed by pandoc snippets
    % extracted from the output of `pandoc -s`
    \providecommand{\tightlist}{%
      \setlength{\itemsep}{0pt}\setlength{\parskip}{0pt}}
    \DefineVerbatimEnvironment{Highlighting}{Verbatim}{commandchars=\\\{\}}
    % Add ',fontsize=\small' for more characters per line
    \newenvironment{Shaded}{}{}
    \newcommand{\KeywordTok}[1]{\textcolor[rgb]{0.00,0.44,0.13}{\textbf{{#1}}}}
    \newcommand{\DataTypeTok}[1]{\textcolor[rgb]{0.56,0.13,0.00}{{#1}}}
    \newcommand{\DecValTok}[1]{\textcolor[rgb]{0.25,0.63,0.44}{{#1}}}
    \newcommand{\BaseNTok}[1]{\textcolor[rgb]{0.25,0.63,0.44}{{#1}}}
    \newcommand{\FloatTok}[1]{\textcolor[rgb]{0.25,0.63,0.44}{{#1}}}
    \newcommand{\CharTok}[1]{\textcolor[rgb]{0.25,0.44,0.63}{{#1}}}
    \newcommand{\StringTok}[1]{\textcolor[rgb]{0.25,0.44,0.63}{{#1}}}
    \newcommand{\CommentTok}[1]{\textcolor[rgb]{0.38,0.63,0.69}{\textit{{#1}}}}
    \newcommand{\OtherTok}[1]{\textcolor[rgb]{0.00,0.44,0.13}{{#1}}}
    \newcommand{\AlertTok}[1]{\textcolor[rgb]{1.00,0.00,0.00}{\textbf{{#1}}}}
    \newcommand{\FunctionTok}[1]{\textcolor[rgb]{0.02,0.16,0.49}{{#1}}}
    \newcommand{\RegionMarkerTok}[1]{{#1}}
    \newcommand{\ErrorTok}[1]{\textcolor[rgb]{1.00,0.00,0.00}{\textbf{{#1}}}}
    \newcommand{\NormalTok}[1]{{#1}}
    
    % Additional commands for more recent versions of Pandoc
    \newcommand{\ConstantTok}[1]{\textcolor[rgb]{0.53,0.00,0.00}{{#1}}}
    \newcommand{\SpecialCharTok}[1]{\textcolor[rgb]{0.25,0.44,0.63}{{#1}}}
    \newcommand{\VerbatimStringTok}[1]{\textcolor[rgb]{0.25,0.44,0.63}{{#1}}}
    \newcommand{\SpecialStringTok}[1]{\textcolor[rgb]{0.73,0.40,0.53}{{#1}}}
    \newcommand{\ImportTok}[1]{{#1}}
    \newcommand{\DocumentationTok}[1]{\textcolor[rgb]{0.73,0.13,0.13}{\textit{{#1}}}}
    \newcommand{\AnnotationTok}[1]{\textcolor[rgb]{0.38,0.63,0.69}{\textbf{\textit{{#1}}}}}
    \newcommand{\CommentVarTok}[1]{\textcolor[rgb]{0.38,0.63,0.69}{\textbf{\textit{{#1}}}}}
    \newcommand{\VariableTok}[1]{\textcolor[rgb]{0.10,0.09,0.49}{{#1}}}
    \newcommand{\ControlFlowTok}[1]{\textcolor[rgb]{0.00,0.44,0.13}{\textbf{{#1}}}}
    \newcommand{\OperatorTok}[1]{\textcolor[rgb]{0.40,0.40,0.40}{{#1}}}
    \newcommand{\BuiltInTok}[1]{{#1}}
    \newcommand{\ExtensionTok}[1]{{#1}}
    \newcommand{\PreprocessorTok}[1]{\textcolor[rgb]{0.74,0.48,0.00}{{#1}}}
    \newcommand{\AttributeTok}[1]{\textcolor[rgb]{0.49,0.56,0.16}{{#1}}}
    \newcommand{\InformationTok}[1]{\textcolor[rgb]{0.38,0.63,0.69}{\textbf{\textit{{#1}}}}}
    \newcommand{\WarningTok}[1]{\textcolor[rgb]{0.38,0.63,0.69}{\textbf{\textit{{#1}}}}}
    
    
    % Define a nice break command that doesn't care if a line doesn't already
    % exist.
    \def\br{\hspace*{\fill} \\* }
    % Math Jax compatibility definitions
    \def\gt{>}
    \def\lt{<}
    \let\Oldtex\TeX
    \let\Oldlatex\LaTeX
    \renewcommand{\TeX}{\textrm{\Oldtex}}
    \renewcommand{\LaTeX}{\textrm{\Oldlatex}}
    % Document parameters
    % Document title
    \title{Notebook}
    
    
    
    
    
% Pygments definitions
\makeatletter
\def\PY@reset{\let\PY@it=\relax \let\PY@bf=\relax%
    \let\PY@ul=\relax \let\PY@tc=\relax%
    \let\PY@bc=\relax \let\PY@ff=\relax}
\def\PY@tok#1{\csname PY@tok@#1\endcsname}
\def\PY@toks#1+{\ifx\relax#1\empty\else%
    \PY@tok{#1}\expandafter\PY@toks\fi}
\def\PY@do#1{\PY@bc{\PY@tc{\PY@ul{%
    \PY@it{\PY@bf{\PY@ff{#1}}}}}}}
\def\PY#1#2{\PY@reset\PY@toks#1+\relax+\PY@do{#2}}

\expandafter\def\csname PY@tok@w\endcsname{\def\PY@tc##1{\textcolor[rgb]{0.73,0.73,0.73}{##1}}}
\expandafter\def\csname PY@tok@c\endcsname{\let\PY@it=\textit\def\PY@tc##1{\textcolor[rgb]{0.25,0.50,0.50}{##1}}}
\expandafter\def\csname PY@tok@cp\endcsname{\def\PY@tc##1{\textcolor[rgb]{0.74,0.48,0.00}{##1}}}
\expandafter\def\csname PY@tok@k\endcsname{\let\PY@bf=\textbf\def\PY@tc##1{\textcolor[rgb]{0.00,0.50,0.00}{##1}}}
\expandafter\def\csname PY@tok@kp\endcsname{\def\PY@tc##1{\textcolor[rgb]{0.00,0.50,0.00}{##1}}}
\expandafter\def\csname PY@tok@kt\endcsname{\def\PY@tc##1{\textcolor[rgb]{0.69,0.00,0.25}{##1}}}
\expandafter\def\csname PY@tok@o\endcsname{\def\PY@tc##1{\textcolor[rgb]{0.40,0.40,0.40}{##1}}}
\expandafter\def\csname PY@tok@ow\endcsname{\let\PY@bf=\textbf\def\PY@tc##1{\textcolor[rgb]{0.67,0.13,1.00}{##1}}}
\expandafter\def\csname PY@tok@nb\endcsname{\def\PY@tc##1{\textcolor[rgb]{0.00,0.50,0.00}{##1}}}
\expandafter\def\csname PY@tok@nf\endcsname{\def\PY@tc##1{\textcolor[rgb]{0.00,0.00,1.00}{##1}}}
\expandafter\def\csname PY@tok@nc\endcsname{\let\PY@bf=\textbf\def\PY@tc##1{\textcolor[rgb]{0.00,0.00,1.00}{##1}}}
\expandafter\def\csname PY@tok@nn\endcsname{\let\PY@bf=\textbf\def\PY@tc##1{\textcolor[rgb]{0.00,0.00,1.00}{##1}}}
\expandafter\def\csname PY@tok@ne\endcsname{\let\PY@bf=\textbf\def\PY@tc##1{\textcolor[rgb]{0.82,0.25,0.23}{##1}}}
\expandafter\def\csname PY@tok@nv\endcsname{\def\PY@tc##1{\textcolor[rgb]{0.10,0.09,0.49}{##1}}}
\expandafter\def\csname PY@tok@no\endcsname{\def\PY@tc##1{\textcolor[rgb]{0.53,0.00,0.00}{##1}}}
\expandafter\def\csname PY@tok@nl\endcsname{\def\PY@tc##1{\textcolor[rgb]{0.63,0.63,0.00}{##1}}}
\expandafter\def\csname PY@tok@ni\endcsname{\let\PY@bf=\textbf\def\PY@tc##1{\textcolor[rgb]{0.60,0.60,0.60}{##1}}}
\expandafter\def\csname PY@tok@na\endcsname{\def\PY@tc##1{\textcolor[rgb]{0.49,0.56,0.16}{##1}}}
\expandafter\def\csname PY@tok@nt\endcsname{\let\PY@bf=\textbf\def\PY@tc##1{\textcolor[rgb]{0.00,0.50,0.00}{##1}}}
\expandafter\def\csname PY@tok@nd\endcsname{\def\PY@tc##1{\textcolor[rgb]{0.67,0.13,1.00}{##1}}}
\expandafter\def\csname PY@tok@s\endcsname{\def\PY@tc##1{\textcolor[rgb]{0.73,0.13,0.13}{##1}}}
\expandafter\def\csname PY@tok@sd\endcsname{\let\PY@it=\textit\def\PY@tc##1{\textcolor[rgb]{0.73,0.13,0.13}{##1}}}
\expandafter\def\csname PY@tok@si\endcsname{\let\PY@bf=\textbf\def\PY@tc##1{\textcolor[rgb]{0.73,0.40,0.53}{##1}}}
\expandafter\def\csname PY@tok@se\endcsname{\let\PY@bf=\textbf\def\PY@tc##1{\textcolor[rgb]{0.73,0.40,0.13}{##1}}}
\expandafter\def\csname PY@tok@sr\endcsname{\def\PY@tc##1{\textcolor[rgb]{0.73,0.40,0.53}{##1}}}
\expandafter\def\csname PY@tok@ss\endcsname{\def\PY@tc##1{\textcolor[rgb]{0.10,0.09,0.49}{##1}}}
\expandafter\def\csname PY@tok@sx\endcsname{\def\PY@tc##1{\textcolor[rgb]{0.00,0.50,0.00}{##1}}}
\expandafter\def\csname PY@tok@m\endcsname{\def\PY@tc##1{\textcolor[rgb]{0.40,0.40,0.40}{##1}}}
\expandafter\def\csname PY@tok@gh\endcsname{\let\PY@bf=\textbf\def\PY@tc##1{\textcolor[rgb]{0.00,0.00,0.50}{##1}}}
\expandafter\def\csname PY@tok@gu\endcsname{\let\PY@bf=\textbf\def\PY@tc##1{\textcolor[rgb]{0.50,0.00,0.50}{##1}}}
\expandafter\def\csname PY@tok@gd\endcsname{\def\PY@tc##1{\textcolor[rgb]{0.63,0.00,0.00}{##1}}}
\expandafter\def\csname PY@tok@gi\endcsname{\def\PY@tc##1{\textcolor[rgb]{0.00,0.63,0.00}{##1}}}
\expandafter\def\csname PY@tok@gr\endcsname{\def\PY@tc##1{\textcolor[rgb]{1.00,0.00,0.00}{##1}}}
\expandafter\def\csname PY@tok@ge\endcsname{\let\PY@it=\textit}
\expandafter\def\csname PY@tok@gs\endcsname{\let\PY@bf=\textbf}
\expandafter\def\csname PY@tok@gp\endcsname{\let\PY@bf=\textbf\def\PY@tc##1{\textcolor[rgb]{0.00,0.00,0.50}{##1}}}
\expandafter\def\csname PY@tok@go\endcsname{\def\PY@tc##1{\textcolor[rgb]{0.53,0.53,0.53}{##1}}}
\expandafter\def\csname PY@tok@gt\endcsname{\def\PY@tc##1{\textcolor[rgb]{0.00,0.27,0.87}{##1}}}
\expandafter\def\csname PY@tok@err\endcsname{\def\PY@bc##1{\setlength{\fboxsep}{0pt}\fcolorbox[rgb]{1.00,0.00,0.00}{1,1,1}{\strut ##1}}}
\expandafter\def\csname PY@tok@kc\endcsname{\let\PY@bf=\textbf\def\PY@tc##1{\textcolor[rgb]{0.00,0.50,0.00}{##1}}}
\expandafter\def\csname PY@tok@kd\endcsname{\let\PY@bf=\textbf\def\PY@tc##1{\textcolor[rgb]{0.00,0.50,0.00}{##1}}}
\expandafter\def\csname PY@tok@kn\endcsname{\let\PY@bf=\textbf\def\PY@tc##1{\textcolor[rgb]{0.00,0.50,0.00}{##1}}}
\expandafter\def\csname PY@tok@kr\endcsname{\let\PY@bf=\textbf\def\PY@tc##1{\textcolor[rgb]{0.00,0.50,0.00}{##1}}}
\expandafter\def\csname PY@tok@bp\endcsname{\def\PY@tc##1{\textcolor[rgb]{0.00,0.50,0.00}{##1}}}
\expandafter\def\csname PY@tok@fm\endcsname{\def\PY@tc##1{\textcolor[rgb]{0.00,0.00,1.00}{##1}}}
\expandafter\def\csname PY@tok@vc\endcsname{\def\PY@tc##1{\textcolor[rgb]{0.10,0.09,0.49}{##1}}}
\expandafter\def\csname PY@tok@vg\endcsname{\def\PY@tc##1{\textcolor[rgb]{0.10,0.09,0.49}{##1}}}
\expandafter\def\csname PY@tok@vi\endcsname{\def\PY@tc##1{\textcolor[rgb]{0.10,0.09,0.49}{##1}}}
\expandafter\def\csname PY@tok@vm\endcsname{\def\PY@tc##1{\textcolor[rgb]{0.10,0.09,0.49}{##1}}}
\expandafter\def\csname PY@tok@sa\endcsname{\def\PY@tc##1{\textcolor[rgb]{0.73,0.13,0.13}{##1}}}
\expandafter\def\csname PY@tok@sb\endcsname{\def\PY@tc##1{\textcolor[rgb]{0.73,0.13,0.13}{##1}}}
\expandafter\def\csname PY@tok@sc\endcsname{\def\PY@tc##1{\textcolor[rgb]{0.73,0.13,0.13}{##1}}}
\expandafter\def\csname PY@tok@dl\endcsname{\def\PY@tc##1{\textcolor[rgb]{0.73,0.13,0.13}{##1}}}
\expandafter\def\csname PY@tok@s2\endcsname{\def\PY@tc##1{\textcolor[rgb]{0.73,0.13,0.13}{##1}}}
\expandafter\def\csname PY@tok@sh\endcsname{\def\PY@tc##1{\textcolor[rgb]{0.73,0.13,0.13}{##1}}}
\expandafter\def\csname PY@tok@s1\endcsname{\def\PY@tc##1{\textcolor[rgb]{0.73,0.13,0.13}{##1}}}
\expandafter\def\csname PY@tok@mb\endcsname{\def\PY@tc##1{\textcolor[rgb]{0.40,0.40,0.40}{##1}}}
\expandafter\def\csname PY@tok@mf\endcsname{\def\PY@tc##1{\textcolor[rgb]{0.40,0.40,0.40}{##1}}}
\expandafter\def\csname PY@tok@mh\endcsname{\def\PY@tc##1{\textcolor[rgb]{0.40,0.40,0.40}{##1}}}
\expandafter\def\csname PY@tok@mi\endcsname{\def\PY@tc##1{\textcolor[rgb]{0.40,0.40,0.40}{##1}}}
\expandafter\def\csname PY@tok@il\endcsname{\def\PY@tc##1{\textcolor[rgb]{0.40,0.40,0.40}{##1}}}
\expandafter\def\csname PY@tok@mo\endcsname{\def\PY@tc##1{\textcolor[rgb]{0.40,0.40,0.40}{##1}}}
\expandafter\def\csname PY@tok@ch\endcsname{\let\PY@it=\textit\def\PY@tc##1{\textcolor[rgb]{0.25,0.50,0.50}{##1}}}
\expandafter\def\csname PY@tok@cm\endcsname{\let\PY@it=\textit\def\PY@tc##1{\textcolor[rgb]{0.25,0.50,0.50}{##1}}}
\expandafter\def\csname PY@tok@cpf\endcsname{\let\PY@it=\textit\def\PY@tc##1{\textcolor[rgb]{0.25,0.50,0.50}{##1}}}
\expandafter\def\csname PY@tok@c1\endcsname{\let\PY@it=\textit\def\PY@tc##1{\textcolor[rgb]{0.25,0.50,0.50}{##1}}}
\expandafter\def\csname PY@tok@cs\endcsname{\let\PY@it=\textit\def\PY@tc##1{\textcolor[rgb]{0.25,0.50,0.50}{##1}}}

\def\PYZbs{\char`\\}
\def\PYZus{\char`\_}
\def\PYZob{\char`\{}
\def\PYZcb{\char`\}}
\def\PYZca{\char`\^}
\def\PYZam{\char`\&}
\def\PYZlt{\char`\<}
\def\PYZgt{\char`\>}
\def\PYZsh{\char`\#}
\def\PYZpc{\char`\%}
\def\PYZdl{\char`\$}
\def\PYZhy{\char`\-}
\def\PYZsq{\char`\'}
\def\PYZdq{\char`\"}
\def\PYZti{\char`\~}
% for compatibility with earlier versions
\def\PYZat{@}
\def\PYZlb{[}
\def\PYZrb{]}
\makeatother


    % For linebreaks inside Verbatim environment from package fancyvrb. 
    \makeatletter
        \newbox\Wrappedcontinuationbox 
        \newbox\Wrappedvisiblespacebox 
        \newcommand*\Wrappedvisiblespace {\textcolor{red}{\textvisiblespace}} 
        \newcommand*\Wrappedcontinuationsymbol {\textcolor{red}{\llap{\tiny$\m@th\hookrightarrow$}}} 
        \newcommand*\Wrappedcontinuationindent {3ex } 
        \newcommand*\Wrappedafterbreak {\kern\Wrappedcontinuationindent\copy\Wrappedcontinuationbox} 
        % Take advantage of the already applied Pygments mark-up to insert 
        % potential linebreaks for TeX processing. 
        %        {, <, #, %, $, ' and ": go to next line. 
        %        _, }, ^, &, >, - and ~: stay at end of broken line. 
        % Use of \textquotesingle for straight quote. 
        \newcommand*\Wrappedbreaksatspecials {% 
            \def\PYGZus{\discretionary{\char`\_}{\Wrappedafterbreak}{\char`\_}}% 
            \def\PYGZob{\discretionary{}{\Wrappedafterbreak\char`\{}{\char`\{}}% 
            \def\PYGZcb{\discretionary{\char`\}}{\Wrappedafterbreak}{\char`\}}}% 
            \def\PYGZca{\discretionary{\char`\^}{\Wrappedafterbreak}{\char`\^}}% 
            \def\PYGZam{\discretionary{\char`\&}{\Wrappedafterbreak}{\char`\&}}% 
            \def\PYGZlt{\discretionary{}{\Wrappedafterbreak\char`\<}{\char`\<}}% 
            \def\PYGZgt{\discretionary{\char`\>}{\Wrappedafterbreak}{\char`\>}}% 
            \def\PYGZsh{\discretionary{}{\Wrappedafterbreak\char`\#}{\char`\#}}% 
            \def\PYGZpc{\discretionary{}{\Wrappedafterbreak\char`\%}{\char`\%}}% 
            \def\PYGZdl{\discretionary{}{\Wrappedafterbreak\char`\$}{\char`\$}}% 
            \def\PYGZhy{\discretionary{\char`\-}{\Wrappedafterbreak}{\char`\-}}% 
            \def\PYGZsq{\discretionary{}{\Wrappedafterbreak\textquotesingle}{\textquotesingle}}% 
            \def\PYGZdq{\discretionary{}{\Wrappedafterbreak\char`\"}{\char`\"}}% 
            \def\PYGZti{\discretionary{\char`\~}{\Wrappedafterbreak}{\char`\~}}% 
        } 
        % Some characters . , ; ? ! / are not pygmentized. 
        % This macro makes them "active" and they will insert potential linebreaks 
        \newcommand*\Wrappedbreaksatpunct {% 
            \lccode`\~`\.\lowercase{\def~}{\discretionary{\hbox{\char`\.}}{\Wrappedafterbreak}{\hbox{\char`\.}}}% 
            \lccode`\~`\,\lowercase{\def~}{\discretionary{\hbox{\char`\,}}{\Wrappedafterbreak}{\hbox{\char`\,}}}% 
            \lccode`\~`\;\lowercase{\def~}{\discretionary{\hbox{\char`\;}}{\Wrappedafterbreak}{\hbox{\char`\;}}}% 
            \lccode`\~`\:\lowercase{\def~}{\discretionary{\hbox{\char`\:}}{\Wrappedafterbreak}{\hbox{\char`\:}}}% 
            \lccode`\~`\?\lowercase{\def~}{\discretionary{\hbox{\char`\?}}{\Wrappedafterbreak}{\hbox{\char`\?}}}% 
            \lccode`\~`\!\lowercase{\def~}{\discretionary{\hbox{\char`\!}}{\Wrappedafterbreak}{\hbox{\char`\!}}}% 
            \lccode`\~`\/\lowercase{\def~}{\discretionary{\hbox{\char`\/}}{\Wrappedafterbreak}{\hbox{\char`\/}}}% 
            \catcode`\.\active
            \catcode`\,\active 
            \catcode`\;\active
            \catcode`\:\active
            \catcode`\?\active
            \catcode`\!\active
            \catcode`\/\active 
            \lccode`\~`\~ 	
        }
    \makeatother

    \let\OriginalVerbatim=\Verbatim
    \makeatletter
    \renewcommand{\Verbatim}[1][1]{%
        %\parskip\z@skip
        \sbox\Wrappedcontinuationbox {\Wrappedcontinuationsymbol}%
        \sbox\Wrappedvisiblespacebox {\FV@SetupFont\Wrappedvisiblespace}%
        \def\FancyVerbFormatLine ##1{\hsize\linewidth
            \vtop{\raggedright\hyphenpenalty\z@\exhyphenpenalty\z@
                \doublehyphendemerits\z@\finalhyphendemerits\z@
                \strut ##1\strut}%
        }%
        % If the linebreak is at a space, the latter will be displayed as visible
        % space at end of first line, and a continuation symbol starts next line.
        % Stretch/shrink are however usually zero for typewriter font.
        \def\FV@Space {%
            \nobreak\hskip\z@ plus\fontdimen3\font minus\fontdimen4\font
            \discretionary{\copy\Wrappedvisiblespacebox}{\Wrappedafterbreak}
            {\kern\fontdimen2\font}%
        }%
        
        % Allow breaks at special characters using \PYG... macros.
        \Wrappedbreaksatspecials
        % Breaks at punctuation characters . , ; ? ! and / need catcode=\active 	
        \OriginalVerbatim[#1,codes*=\Wrappedbreaksatpunct]%
    }
    \makeatother

    % Exact colors from NB
    \definecolor{incolor}{HTML}{303F9F}
    \definecolor{outcolor}{HTML}{D84315}
    \definecolor{cellborder}{HTML}{CFCFCF}
    \definecolor{cellbackground}{HTML}{F7F7F7}
    
    % prompt
    \makeatletter
    \newcommand{\boxspacing}{\kern\kvtcb@left@rule\kern\kvtcb@boxsep}
    \makeatother
    \newcommand{\prompt}[4]{
        \ttfamily\llap{{\color{#2}[#3]:\hspace{3pt}#4}}\vspace{-\baselineskip}
    }
    

    
    % Prevent overflowing lines due to hard-to-break entities
    \sloppy 
    % Setup hyperref package
    \hypersetup{
      breaklinks=true,  % so long urls are correctly broken across lines
      colorlinks=true,
      urlcolor=urlcolor,
      linkcolor=linkcolor,
      citecolor=citecolor,
      }
    % Slightly bigger margins than the latex defaults
    
    \geometry{verbose,tmargin=1in,bmargin=1in,lmargin=1in,rmargin=1in}
    
    

\begin{document}
    
    \maketitle
    
    

    
    \hypertarget{fastpages-notebook-blog-post}{%
\section{Fastpages Notebook Blog
Post}\label{fastpages-notebook-blog-post}}
\label{sec:01-test}
    \begin{quote}
A tutorial of fastpages for Jupyter notebooks.
\end{quote}

\begin{itemize}
\tightlist
\item
  toc: true
\item
  badges: true
\item
  comments: true
\item
  categories: {[}jupyter{]}
\item
  image: images/chart-preview.png
\end{itemize}

    \hypertarget{about}{%
\section{About}\label{about}}

This notebook is a demonstration of some of capabilities of
\href{https://github.com/fastai/fastpages}{fastpages} with notebooks.

With \texttt{fastpages} you can save your jupyter notebooks into the
\texttt{\_notebooks} folder at the root of your repository, and they
will be automatically be converted to Jekyll compliant blog posts!

    \hypertarget{front-matter}{%
\subsection{Front Matter}\label{front-matter}}

The first cell in your Jupyter Notebook or markdown blog post contains
front matter. Front matter is metadata that can turn on/off options in
your Notebook. It is formatted like this:

\begin{verbatim}
# "My Title"
> "Awesome summary"

- toc: true- branch: master- badges: true
- comments: true
- author: Hamel Husain & Jeremy Howard
- categories: [fastpages, jupyter]
\end{verbatim}

\begin{itemize}
\tightlist
\item
  Setting \texttt{toc:\ true} will automatically generate a table of
  contents
\item
  Setting \texttt{badges:\ true} will automatically include GitHub and
  Google Colab links to your notebook.
\item
  Setting \texttt{comments:\ true} will enable commenting on your blog
  post, powered by
  \href{https://github.com/utterance/utterances}{utterances}.
\end{itemize}

The title and description need to be enclosed in double quotes only if
they include special characters such as a colon. More details and
options for front matter can be viewed on the
\href{https://github.com/fastai/fastpages\#front-matter-related-options}{front
matter section} of the README.

    \hypertarget{markdown-shortcuts}{%
\subsection{Markdown Shortcuts}\label{markdown-shortcuts}}

    A \texttt{\#hide} comment at the top of any code cell will hide
\textbf{both the input and output} of that cell in your blog post.

A \texttt{\#hide\_input} comment at the top of any code cell will
\textbf{only hide the input} of that cell.

    \begin{tcolorbox}[breakable, size=fbox, boxrule=1pt, pad at break*=1mm,colback=cellbackground, colframe=cellborder]
\prompt{In}{incolor}{1}{\boxspacing}
\begin{Verbatim}[commandchars=\\\{\}]
\PY{c+c1}{\PYZsh{}hide\PYZus{}input}
\PY{n+nb}{print}\PY{p}{(}\PY{l+s+s1}{\PYZsq{}}\PY{l+s+s1}{The comment \PYZsh{}hide\PYZus{}input was used to hide the code that produced this.}\PY{l+s+s1}{\PYZsq{}}\PY{p}{)}
\end{Verbatim}
\end{tcolorbox}

    \begin{Verbatim}[commandchars=\\\{\}]
The comment \#hide\_input was used to hide the code that produced this.
    \end{Verbatim}

    put a \texttt{\#collapse-hide} flag at the top of any cell if you want
to \textbf{hide} that cell by default, but give the reader the option to
show it:

    \begin{tcolorbox}[breakable, size=fbox, boxrule=1pt, pad at break*=1mm,colback=cellbackground, colframe=cellborder]
\prompt{In}{incolor}{6}{\boxspacing}
\begin{Verbatim}[commandchars=\\\{\}]
\PY{c+c1}{\PYZsh{}collapse\PYZhy{}hide}
\PY{k+kn}{import} \PY{n+nn}{pandas} \PY{k}{as} \PY{n+nn}{pd}
\PY{k+kn}{import} \PY{n+nn}{altair} \PY{k}{as} \PY{n+nn}{alt}
\end{Verbatim}
\end{tcolorbox}

    put a \texttt{\#collapse-show} flag at the top of any cell if you want
to \textbf{show} that cell by default, but give the reader the option to
hide it:

    \begin{tcolorbox}[breakable, size=fbox, boxrule=1pt, pad at break*=1mm,colback=cellbackground, colframe=cellborder]
\prompt{In}{incolor}{4}{\boxspacing}
\begin{Verbatim}[commandchars=\\\{\}]
\PY{c+c1}{\PYZsh{}collapse\PYZhy{}show}
\PY{n}{cars} \PY{o}{=} \PY{l+s+s1}{\PYZsq{}}\PY{l+s+s1}{https://vega.github.io/vega\PYZhy{}datasets/data/cars.json}\PY{l+s+s1}{\PYZsq{}}
\PY{n}{movies} \PY{o}{=} \PY{l+s+s1}{\PYZsq{}}\PY{l+s+s1}{https://vega.github.io/vega\PYZhy{}datasets/data/movies.json}\PY{l+s+s1}{\PYZsq{}}
\PY{n}{sp500} \PY{o}{=} \PY{l+s+s1}{\PYZsq{}}\PY{l+s+s1}{https://vega.github.io/vega\PYZhy{}datasets/data/sp500.csv}\PY{l+s+s1}{\PYZsq{}}
\PY{n}{stocks} \PY{o}{=} \PY{l+s+s1}{\PYZsq{}}\PY{l+s+s1}{https://vega.github.io/vega\PYZhy{}datasets/data/stocks.csv}\PY{l+s+s1}{\PYZsq{}}
\PY{n}{flights} \PY{o}{=} \PY{l+s+s1}{\PYZsq{}}\PY{l+s+s1}{https://vega.github.io/vega\PYZhy{}datasets/data/flights\PYZhy{}5k.json}\PY{l+s+s1}{\PYZsq{}}
\end{Verbatim}
\end{tcolorbox}

    \hypertarget{interactive-charts-with-altair}{%
\subsection{Interactive Charts With
Altair}\label{interactive-charts-with-altair}}

Charts made with Altair remain interactive. Example charts taken from
\href{https://github.com/uwdata/visualization-curriculum}{this repo},
specifically
\href{https://github.com/uwdata/visualization-curriculum/blob/master/altair_interaction.ipynb}{this
notebook}.

    \begin{tcolorbox}[breakable, size=fbox, boxrule=1pt, pad at break*=1mm,colback=cellbackground, colframe=cellborder]
\prompt{In}{incolor}{5}{\boxspacing}
\begin{Verbatim}[commandchars=\\\{\}]
\PY{c+c1}{\PYZsh{} hide}
\PY{n}{df} \PY{o}{=} \PY{n}{pd}\PY{o}{.}\PY{n}{read\PYZus{}json}\PY{p}{(}\PY{n}{movies}\PY{p}{)} \PY{c+c1}{\PYZsh{} load movies data}
\PY{n}{genres} \PY{o}{=} \PY{n}{df}\PY{p}{[}\PY{l+s+s1}{\PYZsq{}}\PY{l+s+s1}{Major\PYZus{}Genre}\PY{l+s+s1}{\PYZsq{}}\PY{p}{]}\PY{o}{.}\PY{n}{unique}\PY{p}{(}\PY{p}{)} \PY{c+c1}{\PYZsh{} get unique field values}
\PY{n}{genres} \PY{o}{=} \PY{n+nb}{list}\PY{p}{(}\PY{n+nb}{filter}\PY{p}{(}\PY{k}{lambda} \PY{n}{d}\PY{p}{:} \PY{n}{d} \PY{o+ow}{is} \PY{o+ow}{not} \PY{k+kc}{None}\PY{p}{,} \PY{n}{genres}\PY{p}{)}\PY{p}{)} \PY{c+c1}{\PYZsh{} filter out None values}
\PY{n}{genres}\PY{o}{.}\PY{n}{sort}\PY{p}{(}\PY{p}{)} \PY{c+c1}{\PYZsh{} sort alphabetically}
\end{Verbatim}
\end{tcolorbox}

    \begin{tcolorbox}[breakable, size=fbox, boxrule=1pt, pad at break*=1mm,colback=cellbackground, colframe=cellborder]
\prompt{In}{incolor}{4}{\boxspacing}
\begin{Verbatim}[commandchars=\\\{\}]
\PY{c+c1}{\PYZsh{}hide}
\PY{n}{mpaa} \PY{o}{=} \PY{p}{[}\PY{l+s+s1}{\PYZsq{}}\PY{l+s+s1}{G}\PY{l+s+s1}{\PYZsq{}}\PY{p}{,} \PY{l+s+s1}{\PYZsq{}}\PY{l+s+s1}{PG}\PY{l+s+s1}{\PYZsq{}}\PY{p}{,} \PY{l+s+s1}{\PYZsq{}}\PY{l+s+s1}{PG\PYZhy{}13}\PY{l+s+s1}{\PYZsq{}}\PY{p}{,} \PY{l+s+s1}{\PYZsq{}}\PY{l+s+s1}{R}\PY{l+s+s1}{\PYZsq{}}\PY{p}{,} \PY{l+s+s1}{\PYZsq{}}\PY{l+s+s1}{NC\PYZhy{}17}\PY{l+s+s1}{\PYZsq{}}\PY{p}{,} \PY{l+s+s1}{\PYZsq{}}\PY{l+s+s1}{Not Rated}\PY{l+s+s1}{\PYZsq{}}\PY{p}{]}
\end{Verbatim}
\end{tcolorbox}

    \hypertarget{example-1-dropdown}{%
\subsubsection{Example 1: DropDown}\label{example-1-dropdown}}

    \begin{tcolorbox}[breakable, size=fbox, boxrule=1pt, pad at break*=1mm,colback=cellbackground, colframe=cellborder]
\prompt{In}{incolor}{5}{\boxspacing}
\begin{Verbatim}[commandchars=\\\{\}]
\PY{c+c1}{\PYZsh{} single\PYZhy{}value selection over [Major\PYZus{}Genre, MPAA\PYZus{}Rating] pairs}
\PY{c+c1}{\PYZsh{} use specific hard\PYZhy{}wired values as the initial selected values}
\PY{n}{selection} \PY{o}{=} \PY{n}{alt}\PY{o}{.}\PY{n}{selection\PYZus{}single}\PY{p}{(}
    \PY{n}{name}\PY{o}{=}\PY{l+s+s1}{\PYZsq{}}\PY{l+s+s1}{Select}\PY{l+s+s1}{\PYZsq{}}\PY{p}{,}
    \PY{n}{fields}\PY{o}{=}\PY{p}{[}\PY{l+s+s1}{\PYZsq{}}\PY{l+s+s1}{Major\PYZus{}Genre}\PY{l+s+s1}{\PYZsq{}}\PY{p}{,} \PY{l+s+s1}{\PYZsq{}}\PY{l+s+s1}{MPAA\PYZus{}Rating}\PY{l+s+s1}{\PYZsq{}}\PY{p}{]}\PY{p}{,}
    \PY{n}{init}\PY{o}{=}\PY{p}{\PYZob{}}\PY{l+s+s1}{\PYZsq{}}\PY{l+s+s1}{Major\PYZus{}Genre}\PY{l+s+s1}{\PYZsq{}}\PY{p}{:} \PY{l+s+s1}{\PYZsq{}}\PY{l+s+s1}{Drama}\PY{l+s+s1}{\PYZsq{}}\PY{p}{,} \PY{l+s+s1}{\PYZsq{}}\PY{l+s+s1}{MPAA\PYZus{}Rating}\PY{l+s+s1}{\PYZsq{}}\PY{p}{:} \PY{l+s+s1}{\PYZsq{}}\PY{l+s+s1}{R}\PY{l+s+s1}{\PYZsq{}}\PY{p}{\PYZcb{}}\PY{p}{,}
    \PY{n}{bind}\PY{o}{=}\PY{p}{\PYZob{}}\PY{l+s+s1}{\PYZsq{}}\PY{l+s+s1}{Major\PYZus{}Genre}\PY{l+s+s1}{\PYZsq{}}\PY{p}{:} \PY{n}{alt}\PY{o}{.}\PY{n}{binding\PYZus{}select}\PY{p}{(}\PY{n}{options}\PY{o}{=}\PY{n}{genres}\PY{p}{)}\PY{p}{,} \PY{l+s+s1}{\PYZsq{}}\PY{l+s+s1}{MPAA\PYZus{}Rating}\PY{l+s+s1}{\PYZsq{}}\PY{p}{:} \PY{n}{alt}\PY{o}{.}\PY{n}{binding\PYZus{}radio}\PY{p}{(}\PY{n}{options}\PY{o}{=}\PY{n}{mpaa}\PY{p}{)}\PY{p}{\PYZcb{}}
\PY{p}{)}
  
\PY{c+c1}{\PYZsh{} scatter plot, modify opacity based on selection}
\PY{n}{alt}\PY{o}{.}\PY{n}{Chart}\PY{p}{(}\PY{n}{movies}\PY{p}{)}\PY{o}{.}\PY{n}{mark\PYZus{}circle}\PY{p}{(}\PY{p}{)}\PY{o}{.}\PY{n}{add\PYZus{}selection}\PY{p}{(}
    \PY{n}{selection}
\PY{p}{)}\PY{o}{.}\PY{n}{encode}\PY{p}{(}
    \PY{n}{x}\PY{o}{=}\PY{l+s+s1}{\PYZsq{}}\PY{l+s+s1}{Rotten\PYZus{}Tomatoes\PYZus{}Rating:Q}\PY{l+s+s1}{\PYZsq{}}\PY{p}{,}
    \PY{n}{y}\PY{o}{=}\PY{l+s+s1}{\PYZsq{}}\PY{l+s+s1}{IMDB\PYZus{}Rating:Q}\PY{l+s+s1}{\PYZsq{}}\PY{p}{,}
    \PY{n}{tooltip}\PY{o}{=}\PY{l+s+s1}{\PYZsq{}}\PY{l+s+s1}{Title:N}\PY{l+s+s1}{\PYZsq{}}\PY{p}{,}
    \PY{n}{opacity}\PY{o}{=}\PY{n}{alt}\PY{o}{.}\PY{n}{condition}\PY{p}{(}\PY{n}{selection}\PY{p}{,} \PY{n}{alt}\PY{o}{.}\PY{n}{value}\PY{p}{(}\PY{l+m+mf}{0.75}\PY{p}{)}\PY{p}{,} \PY{n}{alt}\PY{o}{.}\PY{n}{value}\PY{p}{(}\PY{l+m+mf}{0.05}\PY{p}{)}\PY{p}{)}
\PY{p}{)}
\end{Verbatim}
\end{tcolorbox}

            \begin{tcolorbox}[breakable, size=fbox, boxrule=.5pt, pad at break*=1mm, opacityfill=0]
\prompt{Out}{outcolor}{5}{\boxspacing}
\begin{Verbatim}[commandchars=\\\{\}]
alt.Chart({\ldots})
\end{Verbatim}
\end{tcolorbox}
        
    \hypertarget{example-2-tooltips}{%
\subsubsection{Example 2: Tooltips}\label{example-2-tooltips}}

    \begin{tcolorbox}[breakable, size=fbox, boxrule=1pt, pad at break*=1mm,colback=cellbackground, colframe=cellborder]
\prompt{In}{incolor}{6}{\boxspacing}
\begin{Verbatim}[commandchars=\\\{\}]
\PY{n}{alt}\PY{o}{.}\PY{n}{Chart}\PY{p}{(}\PY{n}{movies}\PY{p}{)}\PY{o}{.}\PY{n}{mark\PYZus{}circle}\PY{p}{(}\PY{p}{)}\PY{o}{.}\PY{n}{add\PYZus{}selection}\PY{p}{(}
    \PY{n}{alt}\PY{o}{.}\PY{n}{selection\PYZus{}interval}\PY{p}{(}\PY{n}{bind}\PY{o}{=}\PY{l+s+s1}{\PYZsq{}}\PY{l+s+s1}{scales}\PY{l+s+s1}{\PYZsq{}}\PY{p}{,} \PY{n}{encodings}\PY{o}{=}\PY{p}{[}\PY{l+s+s1}{\PYZsq{}}\PY{l+s+s1}{x}\PY{l+s+s1}{\PYZsq{}}\PY{p}{]}\PY{p}{)}
\PY{p}{)}\PY{o}{.}\PY{n}{encode}\PY{p}{(}
    \PY{n}{x}\PY{o}{=}\PY{l+s+s1}{\PYZsq{}}\PY{l+s+s1}{Rotten\PYZus{}Tomatoes\PYZus{}Rating:Q}\PY{l+s+s1}{\PYZsq{}}\PY{p}{,}
    \PY{n}{y}\PY{o}{=}\PY{n}{alt}\PY{o}{.}\PY{n}{Y}\PY{p}{(}\PY{l+s+s1}{\PYZsq{}}\PY{l+s+s1}{IMDB\PYZus{}Rating:Q}\PY{l+s+s1}{\PYZsq{}}\PY{p}{,} \PY{n}{axis}\PY{o}{=}\PY{n}{alt}\PY{o}{.}\PY{n}{Axis}\PY{p}{(}\PY{n}{minExtent}\PY{o}{=}\PY{l+m+mi}{30}\PY{p}{)}\PY{p}{)}\PY{p}{,} \PY{c+c1}{\PYZsh{} use min extent to stabilize axis title placement}
    \PY{n}{tooltip}\PY{o}{=}\PY{p}{[}\PY{l+s+s1}{\PYZsq{}}\PY{l+s+s1}{Title:N}\PY{l+s+s1}{\PYZsq{}}\PY{p}{,} \PY{l+s+s1}{\PYZsq{}}\PY{l+s+s1}{Release\PYZus{}Date:N}\PY{l+s+s1}{\PYZsq{}}\PY{p}{,} \PY{l+s+s1}{\PYZsq{}}\PY{l+s+s1}{IMDB\PYZus{}Rating:Q}\PY{l+s+s1}{\PYZsq{}}\PY{p}{,} \PY{l+s+s1}{\PYZsq{}}\PY{l+s+s1}{Rotten\PYZus{}Tomatoes\PYZus{}Rating:Q}\PY{l+s+s1}{\PYZsq{}}\PY{p}{]}
\PY{p}{)}\PY{o}{.}\PY{n}{properties}\PY{p}{(}
    \PY{n}{width}\PY{o}{=}\PY{l+m+mi}{600}\PY{p}{,}
    \PY{n}{height}\PY{o}{=}\PY{l+m+mi}{400}
\PY{p}{)}
\end{Verbatim}
\end{tcolorbox}

            \begin{tcolorbox}[breakable, size=fbox, boxrule=.5pt, pad at break*=1mm, opacityfill=0]
\prompt{Out}{outcolor}{6}{\boxspacing}
\begin{Verbatim}[commandchars=\\\{\}]
alt.Chart({\ldots})
\end{Verbatim}
\end{tcolorbox}
        
    \hypertarget{example-3-more-tooltips}{%
\subsubsection{Example 3: More Tooltips}\label{example-3-more-tooltips}}

    \begin{tcolorbox}[breakable, size=fbox, boxrule=1pt, pad at break*=1mm,colback=cellbackground, colframe=cellborder]
\prompt{In}{incolor}{7}{\boxspacing}
\begin{Verbatim}[commandchars=\\\{\}]
\PY{c+c1}{\PYZsh{} select a point for which to provide details\PYZhy{}on\PYZhy{}demand}
\PY{n}{label} \PY{o}{=} \PY{n}{alt}\PY{o}{.}\PY{n}{selection\PYZus{}single}\PY{p}{(}
    \PY{n}{encodings}\PY{o}{=}\PY{p}{[}\PY{l+s+s1}{\PYZsq{}}\PY{l+s+s1}{x}\PY{l+s+s1}{\PYZsq{}}\PY{p}{]}\PY{p}{,} \PY{c+c1}{\PYZsh{} limit selection to x\PYZhy{}axis value}
    \PY{n}{on}\PY{o}{=}\PY{l+s+s1}{\PYZsq{}}\PY{l+s+s1}{mouseover}\PY{l+s+s1}{\PYZsq{}}\PY{p}{,}  \PY{c+c1}{\PYZsh{} select on mouseover events}
    \PY{n}{nearest}\PY{o}{=}\PY{k+kc}{True}\PY{p}{,}    \PY{c+c1}{\PYZsh{} select data point nearest the cursor}
    \PY{n}{empty}\PY{o}{=}\PY{l+s+s1}{\PYZsq{}}\PY{l+s+s1}{none}\PY{l+s+s1}{\PYZsq{}}     \PY{c+c1}{\PYZsh{} empty selection includes no data points}
\PY{p}{)}

\PY{c+c1}{\PYZsh{} define our base line chart of stock prices}
\PY{n}{base} \PY{o}{=} \PY{n}{alt}\PY{o}{.}\PY{n}{Chart}\PY{p}{(}\PY{p}{)}\PY{o}{.}\PY{n}{mark\PYZus{}line}\PY{p}{(}\PY{p}{)}\PY{o}{.}\PY{n}{encode}\PY{p}{(}
    \PY{n}{alt}\PY{o}{.}\PY{n}{X}\PY{p}{(}\PY{l+s+s1}{\PYZsq{}}\PY{l+s+s1}{date:T}\PY{l+s+s1}{\PYZsq{}}\PY{p}{)}\PY{p}{,}
    \PY{n}{alt}\PY{o}{.}\PY{n}{Y}\PY{p}{(}\PY{l+s+s1}{\PYZsq{}}\PY{l+s+s1}{price:Q}\PY{l+s+s1}{\PYZsq{}}\PY{p}{,} \PY{n}{scale}\PY{o}{=}\PY{n}{alt}\PY{o}{.}\PY{n}{Scale}\PY{p}{(}\PY{n+nb}{type}\PY{o}{=}\PY{l+s+s1}{\PYZsq{}}\PY{l+s+s1}{log}\PY{l+s+s1}{\PYZsq{}}\PY{p}{)}\PY{p}{)}\PY{p}{,}
    \PY{n}{alt}\PY{o}{.}\PY{n}{Color}\PY{p}{(}\PY{l+s+s1}{\PYZsq{}}\PY{l+s+s1}{symbol:N}\PY{l+s+s1}{\PYZsq{}}\PY{p}{)}
\PY{p}{)}

\PY{n}{alt}\PY{o}{.}\PY{n}{layer}\PY{p}{(}
    \PY{n}{base}\PY{p}{,} \PY{c+c1}{\PYZsh{} base line chart}
    
    \PY{c+c1}{\PYZsh{} add a rule mark to serve as a guide line}
    \PY{n}{alt}\PY{o}{.}\PY{n}{Chart}\PY{p}{(}\PY{p}{)}\PY{o}{.}\PY{n}{mark\PYZus{}rule}\PY{p}{(}\PY{n}{color}\PY{o}{=}\PY{l+s+s1}{\PYZsq{}}\PY{l+s+s1}{\PYZsh{}aaa}\PY{l+s+s1}{\PYZsq{}}\PY{p}{)}\PY{o}{.}\PY{n}{encode}\PY{p}{(}
        \PY{n}{x}\PY{o}{=}\PY{l+s+s1}{\PYZsq{}}\PY{l+s+s1}{date:T}\PY{l+s+s1}{\PYZsq{}}
    \PY{p}{)}\PY{o}{.}\PY{n}{transform\PYZus{}filter}\PY{p}{(}\PY{n}{label}\PY{p}{)}\PY{p}{,}
    
    \PY{c+c1}{\PYZsh{} add circle marks for selected time points, hide unselected points}
    \PY{n}{base}\PY{o}{.}\PY{n}{mark\PYZus{}circle}\PY{p}{(}\PY{p}{)}\PY{o}{.}\PY{n}{encode}\PY{p}{(}
        \PY{n}{opacity}\PY{o}{=}\PY{n}{alt}\PY{o}{.}\PY{n}{condition}\PY{p}{(}\PY{n}{label}\PY{p}{,} \PY{n}{alt}\PY{o}{.}\PY{n}{value}\PY{p}{(}\PY{l+m+mi}{1}\PY{p}{)}\PY{p}{,} \PY{n}{alt}\PY{o}{.}\PY{n}{value}\PY{p}{(}\PY{l+m+mi}{0}\PY{p}{)}\PY{p}{)}
    \PY{p}{)}\PY{o}{.}\PY{n}{add\PYZus{}selection}\PY{p}{(}\PY{n}{label}\PY{p}{)}\PY{p}{,}

    \PY{c+c1}{\PYZsh{} add white stroked text to provide a legible background for labels}
    \PY{n}{base}\PY{o}{.}\PY{n}{mark\PYZus{}text}\PY{p}{(}\PY{n}{align}\PY{o}{=}\PY{l+s+s1}{\PYZsq{}}\PY{l+s+s1}{left}\PY{l+s+s1}{\PYZsq{}}\PY{p}{,} \PY{n}{dx}\PY{o}{=}\PY{l+m+mi}{5}\PY{p}{,} \PY{n}{dy}\PY{o}{=}\PY{o}{\PYZhy{}}\PY{l+m+mi}{5}\PY{p}{,} \PY{n}{stroke}\PY{o}{=}\PY{l+s+s1}{\PYZsq{}}\PY{l+s+s1}{white}\PY{l+s+s1}{\PYZsq{}}\PY{p}{,} \PY{n}{strokeWidth}\PY{o}{=}\PY{l+m+mi}{2}\PY{p}{)}\PY{o}{.}\PY{n}{encode}\PY{p}{(}
        \PY{n}{text}\PY{o}{=}\PY{l+s+s1}{\PYZsq{}}\PY{l+s+s1}{price:Q}\PY{l+s+s1}{\PYZsq{}}
    \PY{p}{)}\PY{o}{.}\PY{n}{transform\PYZus{}filter}\PY{p}{(}\PY{n}{label}\PY{p}{)}\PY{p}{,}

    \PY{c+c1}{\PYZsh{} add text labels for stock prices}
    \PY{n}{base}\PY{o}{.}\PY{n}{mark\PYZus{}text}\PY{p}{(}\PY{n}{align}\PY{o}{=}\PY{l+s+s1}{\PYZsq{}}\PY{l+s+s1}{left}\PY{l+s+s1}{\PYZsq{}}\PY{p}{,} \PY{n}{dx}\PY{o}{=}\PY{l+m+mi}{5}\PY{p}{,} \PY{n}{dy}\PY{o}{=}\PY{o}{\PYZhy{}}\PY{l+m+mi}{5}\PY{p}{)}\PY{o}{.}\PY{n}{encode}\PY{p}{(}
        \PY{n}{text}\PY{o}{=}\PY{l+s+s1}{\PYZsq{}}\PY{l+s+s1}{price:Q}\PY{l+s+s1}{\PYZsq{}}
    \PY{p}{)}\PY{o}{.}\PY{n}{transform\PYZus{}filter}\PY{p}{(}\PY{n}{label}\PY{p}{)}\PY{p}{,}
    
    \PY{n}{data}\PY{o}{=}\PY{n}{stocks}
\PY{p}{)}\PY{o}{.}\PY{n}{properties}\PY{p}{(}
    \PY{n}{width}\PY{o}{=}\PY{l+m+mi}{700}\PY{p}{,}
    \PY{n}{height}\PY{o}{=}\PY{l+m+mi}{400}
\PY{p}{)}
\end{Verbatim}
\end{tcolorbox}

            \begin{tcolorbox}[breakable, size=fbox, boxrule=.5pt, pad at break*=1mm, opacityfill=0]
\prompt{Out}{outcolor}{7}{\boxspacing}
\begin{Verbatim}[commandchars=\\\{\}]
alt.LayerChart({\ldots})
\end{Verbatim}
\end{tcolorbox}
        
    \hypertarget{data-tables}{%
\subsection{Data Tables}\label{data-tables}}

You can display tables per the usual way in your blog:

    \begin{tcolorbox}[breakable, size=fbox, boxrule=1pt, pad at break*=1mm,colback=cellbackground, colframe=cellborder]
\prompt{In}{incolor}{11}{\boxspacing}
\begin{Verbatim}[commandchars=\\\{\}]
\PY{n}{movies} \PY{o}{=} \PY{l+s+s1}{\PYZsq{}}\PY{l+s+s1}{https://vega.github.io/vega\PYZhy{}datasets/data/movies.json}\PY{l+s+s1}{\PYZsq{}}
\PY{n}{df} \PY{o}{=} \PY{n}{pd}\PY{o}{.}\PY{n}{read\PYZus{}json}\PY{p}{(}\PY{n}{movies}\PY{p}{)}
\PY{c+c1}{\PYZsh{} display table with pandas}
\PY{n}{df}\PY{p}{[}\PY{p}{[}\PY{l+s+s1}{\PYZsq{}}\PY{l+s+s1}{Title}\PY{l+s+s1}{\PYZsq{}}\PY{p}{,} \PY{l+s+s1}{\PYZsq{}}\PY{l+s+s1}{Worldwide\PYZus{}Gross}\PY{l+s+s1}{\PYZsq{}}\PY{p}{,} 
    \PY{l+s+s1}{\PYZsq{}}\PY{l+s+s1}{Production\PYZus{}Budget}\PY{l+s+s1}{\PYZsq{}}\PY{p}{,} \PY{l+s+s1}{\PYZsq{}}\PY{l+s+s1}{Distributor}\PY{l+s+s1}{\PYZsq{}}\PY{p}{,} \PY{l+s+s1}{\PYZsq{}}\PY{l+s+s1}{MPAA\PYZus{}Rating}\PY{l+s+s1}{\PYZsq{}}\PY{p}{,} \PY{l+s+s1}{\PYZsq{}}\PY{l+s+s1}{IMDB\PYZus{}Rating}\PY{l+s+s1}{\PYZsq{}}\PY{p}{,} \PY{l+s+s1}{\PYZsq{}}\PY{l+s+s1}{Rotten\PYZus{}Tomatoes\PYZus{}Rating}\PY{l+s+s1}{\PYZsq{}}\PY{p}{]}\PY{p}{]}\PY{o}{.}\PY{n}{head}\PY{p}{(}\PY{p}{)}
\end{Verbatim}
\end{tcolorbox}

            \begin{tcolorbox}[breakable, size=fbox, boxrule=.5pt, pad at break*=1mm, opacityfill=0]
\prompt{Out}{outcolor}{11}{\boxspacing}
\begin{Verbatim}[commandchars=\\\{\}]
                        Title  Worldwide\_Gross  Production\_Budget Distributor  \textbackslash{}
0              The Land Girls         146083.0          8000000.0    Gramercy
1      First Love, Last Rites          10876.0           300000.0      Strand
2  I Married a Strange Person         203134.0           250000.0   Lionsgate
3        Let's Talk About Sex         373615.0           300000.0   Fine Line
4                        Slam        1087521.0          1000000.0     Trimark

  MPAA\_Rating  IMDB\_Rating  Rotten\_Tomatoes\_Rating
0           R          6.1                     NaN
1           R          6.9                     NaN
2        None          6.8                     NaN
3        None          NaN                    13.0
4           R          3.4                    62.0
\end{Verbatim}
\end{tcolorbox}
        
    \hypertarget{images}{%
\subsection{Images}\label{images}}

\hypertarget{local-images}{%
\subsubsection{Local Images}\label{local-images}}

You can reference local images and they will be copied and rendered on
your blog automatically. You can include these with the following
markdown syntax:

\texttt{!{[}{]}(my\_icons/fastai\_logo.png)}

    \includegraphics{my_icons/fastai_logo.png}

    \hypertarget{remote-images}{%
\subsubsection{Remote Images}\label{remote-images}}

Remote images can be included with the following markdown syntax:

    \hypertarget{animated-gifs}{%
\subsubsection{Animated Gifs}\label{animated-gifs}}

Animated Gifs work, too!

\texttt{!{[}{]}(https://upload.wikimedia.org/wikipedia/commons/7/71/ChessPawnSpecialMoves.gif)}

    \hypertarget{other-elements}{%
\section{Other Elements}\label{other-elements}}

    \hypertarget{github-flavored-emojis}{%
\subsection{GitHub Flavored Emojis}\label{github-flavored-emojis}}

Typing \texttt{I\ give\ this\ post\ two\ :+1:!} will render this:

I give this post two :+1:!

    \hypertarget{tweetcards}{%
\subsection{Tweetcards}\label{tweetcards}}

Typing
\texttt{\textgreater{}\ twitter:\ https://twitter.com/jakevdp/status/1204765621767901185?s=20}
will render this:

\begin{quote}
twitter: https://twitter.com/jakevdp/status/1204765621767901185?s=20
\end{quote}

    \hypertarget{youtube-videos}{%
\subsection{Youtube Videos}\label{youtube-videos}}

Typing \texttt{\textgreater{}\ youtube:\ https://youtu.be/XfoYk\_Z5AkI}
will render this:

\begin{quote}
youtube: https://youtu.be/XfoYk\_Z5AkI
\end{quote}

    \hypertarget{boxes-callouts}{%
\subsection{Boxes / Callouts}\label{boxes-callouts}}

Typing
\texttt{\textgreater{}\ Warning:\ There\ will\ be\ no\ second\ warning!}
will render this:

\begin{quote}
Warning: There will be no second warning!
\end{quote}

Typing
\texttt{\textgreater{}\ Important:\ Pay\ attention!\ It\textquotesingle{}s\ important.}
will render this:

\begin{quote}
Important: Pay attention! It's important.
\end{quote}

Typing \texttt{\textgreater{}\ Tip:\ This\ is\ my\ tip.} will render
this:

\begin{quote}
Tip: This is my tip.
\end{quote}

Typing \texttt{\textgreater{}\ Note:\ Take\ note\ of\ this.} will render
this:

\begin{quote}
Note: Take note of this.
\end{quote}

Typing
\texttt{\textgreater{}\ Note:\ A\ doc\ link\ to\ {[}an\ example\ website:\ fast.ai{]}(https://www.fast.ai/)\ should\ also\ work\ fine.}
will render in the docs:

\begin{quote}
Note: A doc link to \href{https://www.fast.ai/}{an example website:
fast.ai} should also work fine.
\end{quote}

    \hypertarget{footnotes}{%
\subsection{Footnotes}\label{footnotes}}

You can have footnotes in notebooks, however the syntax is different
compared to markdown documents.
\href{https://github.com/fastai/fastpages/blob/master/_fastpages_docs/NOTEBOOK_FOOTNOTES.md}{This
guide provides more detail about this syntax}, which looks like this:

\begin{verbatim}
{% raw %}For example, here is a footnote {% fn 1 %}.
And another {% fn 2 %}
{{ 'This is the footnote.' | fndetail: 1 }}
{{ 'This is the other footnote. You can even have a [link](www.github.com)!' | fndetail: 2 }}{% endraw %}
\end{verbatim}

For example, here is a footnote \{\% fn 1 \%\}.

And another \{\% fn 2 \%\}

\{\{ `This is the footnote.' \textbar{} fndetail: 1 \}\} \{\{ `This is
the other footnote. You can even have a \href{www.github.com}{link}!'
\textbar{} fndetail: 2 \}\}

    \hypertarget{fastpages-notebook-blog-post-2}{%
\section{Fastpages Notebook Blog Post
2}\label{fastpages-notebook-blog-post-2}}
\label{sec:02-ttest}
    \begin{quote}
A tutorial of fastpages for Jupyter notebooks.
\end{quote}

\begin{itemize}
\tightlist
\item
  toc: true
\item
  badges: true
\item
  comments: true
\item
  categories: {[}jupyter{]}
\item
  image: images/chart-preview.png
\end{itemize}


    % Add a bibliography block to the postdoc
    
    
    
\end{document}
